\documentclass[10pt,a4paper]{amsart}
\usepackage[margin=19mm]{geometry}
\usepackage{amsmath,amssymb,mathtools}
\usepackage[scr=rsfs,cal=euler]{mathalfa}
\usepackage{xfrac}
\usepackage[unicode,hidelinks,bookmarks=false]{hyperref}
\newcommand{\ie}{i.e.\ }
\newcommand{\cf}{cf.\ }
\newcommand{\resp}{respectively\ }
\newcommand{\Subsection}{\subsection}
\providecommand{\nobreakdash}{\nobreak\hbox{-}}
\setlength{\emergencystretch}{2em}
\multlinegap0pt
\usepackage{amssymb}
\usepackage{mathtools}
\usepackage[shortlabels]{enumitem}
\newtheorem{thm}{Theorem}
\title{Canonical Truth}
\author{Merlin Carl, Philipp Schlicht}
\date{Source: arXiv:1712.02566v7 (CC BY 4.0)}
\begin{document}
\maketitle

\noindent\emph{Source and licence.} Canonical Truth. Version arXiv:1712.02566v7 (CC BY 4.0), distributed by arXiv under the Creative Commons Attribution 4.0 International licence, http://creativecommons.org/licenses/by/4.0/, as recorded in the arXiv licence metadata for this exact version and shown on the abstract page https://arxiv.org/abs/1712.02566v7. Full licence notice: \texttt{licenses/arxiv-1712.02566-CC-BY-4.0.txt}.\par\smallskip

\noindent\textit{Editorial scope.} Counted unit: the article's thm closure (source0.tex lines 692--698). The article's complete original proof of it is delivered with it (source0.tex lines 699--719, one \texttt{proof} environment of 21 lines). No mathematics is written, repaired or re-derived: the statement and the proof are the article's own lines, in the article's own notation, and no retained line is substituted or re-flowed.  No further source-local context is needed: every cross-reference of every retained line resolves inside the two delivered spans.  Nothing was omitted from any retained span, so the package carries no omission record.  No retained line contains a citation, so the wrapper carries no bibliography and no bibliography entry.  No retained line contains a figure environment, an image-inclusion command or drawing code, so no image is delivered and the package declares no asset dependency.  Editorial formatting only, in addition: an AMS \texttt{amsart} preamble that restates the article's own declarations and macros used by the retained lines, namely the environments \texttt{thm} on separate counters, together with the standard packages those lines need.  The retained environments print in this document as Theorem 1 in order of appearance, whereas the article prints the same items with the numbering of its own full document.  The complete native source of the article as distributed in the arXiv e-print for this exact version is bundled unchanged as \texttt{sources/101266/source0.tex}.
\begin{thm}{\label{sharp closure}}
(i) Let $\phi^{\sharp}$ be the statement `For every $x\subseteq\omega$, $x^{\sharp}$ exists'.
If $\phi^{\sharp}$ holds in $V$, then $\phi^{\sharp}$ is canonically possible.

(ii) If $\phi^{\dagger}$ is the statement that $x^{\dagger}$ exists for  every $x\subseteq\omega$,
and $\phi^{\dagger}$ holds in $V$, then $\phi^{\dagger}$ is canonically possible.
\end{thm}

\begin{proof}
(i) Let us add the unary function symbol $\sharp$ to the language of set theory, with the obvious intended meaning.
For a set $X$, let $\text{Def}^{\sharp}(X)$ be the set of
subsets of $X$ that are definable over $X$ using the language of set theory extended by $\sharp$. 
Now define a class $\hat{L}\subseteq V$ in analogy with $L$ as follows: 

\begin{itemize}
 \item $\hat{L}_{0}=\emptyset$
 \item $\hat{L}_{\alpha+1}=\text{Def}^{\sharp}(\hat{L}_{\alpha})$
 \item $\hat{L}_{\lambda}=\bigcup_{\iota<\lambda}\hat{L}_{\iota}$ for $\iota<\lambda$
 \item Finally, $\hat{L}:=\bigcup_{\iota\in\text{On}}\hat{L}_{\iota}$.
\end{itemize}

By $\Pi^{1}_{2}$-absoluteness of sharps, the sharp function and thus $\hat{L}$ is a definable subclass of $V$ via the recursion
just stated.
It is clear from the definition that $\hat{L}\models\phi^{\sharp}$. That $\hat{L}\models\text{ZFC}$ can be checked similarly
to the fact that ZFC holds in $L$. 
Thus $\hat{L}$ is a canonical model in which $\phi^{\sharp}$ holds.

(ii) is now proved similarly, as the $\dagger$-operation is still $\Pi^{1}_{2}$.
\end{proof}

\end{document}
